\documentclass[aspectratio=169,11pt]{beamer}

\usepackage[T1]{fontenc}
\usepackage[utf8]{inputenc}
\usepackage[italian]{babel}
\usepackage{lmodern}
\usepackage{microtype}
\usepackage{booktabs}

\definecolor{Ink}{HTML}{18212F}
\definecolor{Accent}{HTML}{176B87}
\definecolor{AccentDark}{HTML}{0E4C63}
\definecolor{AccentSoft}{HTML}{E9F3F6}
\definecolor{Muted}{HTML}{596475}

\setbeamercolor{normal text}{fg=Ink,bg=white}
\setbeamercolor{structure}{fg=Accent}
\setbeamercolor{frametitle}{fg=white,bg=AccentDark}
\setbeamercolor{title}{fg=white}
\setbeamercolor{subtitle}{fg=AccentSoft}
\setbeamercolor{block title}{fg=white,bg=AccentDark}
\setbeamercolor{block body}{fg=Ink,bg=AccentSoft}
\setbeamercolor{alerted text}{fg=AccentDark}
\setbeamertemplate{navigation symbols}{}
\setbeamertemplate{itemize item}{\color{Accent}\small$\blacktriangleright$}
\setbeamertemplate{itemize subitem}{\color{Accent}\scriptsize$\bullet$}
\setbeamertemplate{frametitle}{%
  \nointerlineskip
  \begin{beamercolorbox}[wd=\paperwidth,ht=0.9cm,dp=0.25cm,leftskip=0.55cm]{frametitle}
    \usebeamerfont{frametitle}\insertframetitle
  \end{beamercolorbox}
}
\setbeamertemplate{footline}{%
  \leavevmode
  \hbox{%
    \begin{beamercolorbox}[wd=.82\paperwidth,ht=2.6ex,dp=1.1ex,leftskip=.5cm]{author in head/foot}%
      \color{Muted}\scriptsize Francesco Vicidomini · Software Engineer
    \end{beamercolorbox}%
    \begin{beamercolorbox}[wd=.18\paperwidth,ht=2.6ex,dp=1.1ex,rightskip=.5cm plus1fil]{date in head/foot}%
      \color{Muted}\scriptsize \insertframenumber/\inserttotalframenumber
    \end{beamercolorbox}%
  }
}

\title{Francesco Vicidomini}
\subtitle{Software Engineer}
\author{Candidatura per il ruolo Full-stack Developer presso Notarify}
\date{}

\begin{document}

{
  \setbeamercolor{background canvas}{bg=Ink}
  \begin{frame}[plain]
    \vfill
    {\usebeamerfont{title}\color{white}\LARGE\bfseries\inserttitle\par}
    \vspace{0.35cm}
    {\color{AccentSoft}\Large\insertsubtitle\par}
    \vspace{0.75cm}
    {\color{white}\normalsize\insertauthor\par}
    \vfill
    {\color{AccentSoft}\small
      Salerno, Italia \quad·\quad
      \href{mailto:fr.vicidomini94@gmail.com}{fr.vicidomini94@gmail.com}
    }
  \end{frame}
}

\begin{frame}{Il mio profilo}
  \begin{columns}[T,onlytextwidth]
    \begin{column}{0.62\textwidth}
      \begin{itemize}
        \item Oltre \alert{6 anni di esperienza professionale} nello sviluppo di applicazioni web complesse.
        \item Autonomia su analisi, scelte tecniche, integrazione con API, debugging e code review.
        \item Esperienza in cybersecurity, ricerca europea e servizi per la Pubblica Amministrazione.
        \item Coordinamento e mentoring di piccoli team; collaborazione in inglese con partner internazionali.
      \end{itemize}
    \end{column}
    \begin{column}{0.34\textwidth}
      \begin{block}{Posizionamento}
        La mia esperienza principale è su TypeScript e Angular. Conosco Python, l'ho usato spesso per creare degli script per automatizzare dei processi e durante la mia tesi Magistrale.
      \end{block}
    \end{column}
  \end{columns}
\end{frame}

\begin{frame}{Esperienze lavorative}
  \begin{columns}[T,onlytextwidth]
    \begin{column}{0.48\textwidth}
      \begin{block}{Engineering · 2023--oggi}
        \begin{itemize}
          \item Architettura e sviluppo di piattaforme per privacy, cybersecurity e Digital Twin.
          \item REST, WebSocket, MongoDB, Docker, CI/CD e servizi integrazione verso API Spring Boot.
          \item Analisi di colli di bottiglia e riduzione di calcoli server ridondanti.
          \item Coordinamento, review e confronto con stakeholder internazionali.
        \end{itemize}
      \end{block}
    \end{column}
    \begin{column}{0.48\textwidth}
      \begin{block}{NTT DATA · 2020--2023}
        \begin{itemize}
          \item Piattaforma su larga scala per processi amministrativi provinciali.
          \item Sviluppo Angular, API Spring Boot e MySQL.
          \item Modernizzazione di codice legacy e componenti riutilizzabili.
          \item Coordinamento di tre sviluppatori in Agile Scrum.
        \end{itemize}
      \end{block}
    \end{column}
  \end{columns}
\end{frame}

\begin{frame}{Il mio approccio al problem solving}
  \begin{enumerate}
    \item \textbf{Capire il dominio} — chiarisco obiettivi, vincoli e casi limite prima di scegliere la tecnologia.
    \item \textbf{Ridurre il problema} — separo dati, regole di business, integrazioni e interfaccia in parti verificabili.
    \item \textbf{Costruire feedback rapidi} — uso test, logging e piccoli esperimenti per validare le ipotesi.
    \item \textbf{Rendere la soluzione mantenibile} — preferisco API chiare, componenti coesi, testati, revisionati e supportati da una documentazione essenziale.
  \end{enumerate}

  \vspace{0.35cm}
  \begin{block}{Esempio}
    Davanti a risposte API lente, ho seguito il flusso dei dati fino a individuare calcoli ripetuti, poi ho guidato una modifica che li eliminasse senza alterare il risultato atteso.
  \end{block}
\end{frame}

\begin{frame}{Esperienza con Python}
  \begin{columns}[T,onlytextwidth]
    \begin{column}{0.58\textwidth}
      \textbf{Knoxly — tesi magistrale}
      \begin{itemize}
        \item Sviluppo di un modello di ML per la privacy integrato con Telegram Web e Gmail.
        \item Pipeline in due fasi per classificare la sensibilità dei testi.
        \item Sentence embedding multilingue e classificatori specializzati per argomento.
        \item Modello esposto tramite una \alert{REST API in Python e Flask}.
      \end{itemize}
    \end{column}
    \begin{column}{0.38\textwidth}
      \begin{block}{Fondamenta trasferibili}
        \begin{itemize}
          \item OOP e algoritmi
          \item progettazione di API
          \item database relazionali e documentali
          \item test automatici
          \item Docker e CI/CD
          \item lettura e debugging di sistemi back-end
        \end{itemize}
      \end{block}
    \end{column}
  \end{columns}
\end{frame}

\begin{frame}{Esperienze di sviluppo in autonomia}
  Quando una tecnologia mi interessa cerco di studiarla in autonomia e la metto in pratica con esercizi e piccoli che pubblico ul mio \href{https://github.com/CiccioTecchio}{\textbf{GitHub}}.
  I progetti più recenti sono:

  \vspace{0.35cm}
  \begin{columns}[T,onlytextwidth]
    \begin{column}{0.48\textwidth}
      \begin{block}{MyFiscalCode}
        Calcolatore del codice fiscale italiano sviluppato per approfondire
        Next.js, internazionalizzazione, persistenza dei dati e test automatici.

        \vspace{0.25cm}
        \href{https://my-fiscal-code.vercel.app/}{\textbf{App online}}
        \quad·\quad
        \href{https://github.com/CiccioTecchio/MyFiscalCode}{\textbf{Repository GitHub}}
      \end{block}
    \end{column}
    \begin{column}{0.48\textwidth}
      \begin{block}{FrancescoVicidomini.dev}
        Sito personale usato per sperimentare l'evoluzione dello stack front-end:
        nato in Angular, poi riscritto in React e infine in Next.js con deploy su
        Vercel.

        \vspace{0.25cm}
        \href{https://francescovicidomini.dev/}{\textbf{App online}}
        \quad·\quad
        \href{https://github.com/CiccioTecchio/FrancescoVicidomini.dev}{\textbf{Repository GitHub}}
      \end{block}
    \end{column}
  \end{columns}
\end{frame}

\begin{frame}{Perché posso portare valore nel ruolo}
  \begin{columns}[T,onlytextwidth]
    \begin{column}{0.48\textwidth}
      \textbf{Valore immediato}
      \begin{itemize}
        \item Metodo solido di analisi e debugging.
        \item Esperienza reale su API, dati, test e delivery.
        \item Capacità di muovermi in codebase e domini non familiari.
        \item Comunicazione tecnica e collaborazione con il team.
      \end{itemize}
    \end{column}
    \begin{column}{0.48\textwidth}
      \textbf{Percorso su Python}
      \begin{itemize}
        \item Consolidare idiomi, typing e packaging dell'ecosistema.
        \item Applicare da subito test, separazione delle responsabilità e API design.
        \item Imparare sul prodotto reale attraverso feedback e code review.
        \item Trasferire l'esperienza maturata su sistemi web complessi.
      \end{itemize}
    \end{column}
  \end{columns}
\end{frame}

{
  \setbeamercolor{background canvas}{bg=Ink}
  \begin{frame}[plain]
    \vfill
    {\color{white}\LARGE\bfseries Grazie per il vostro tempo\par}
    \vspace{0.45cm}
    {\color{white}
      \href{mailto:fr.vicidomini94@gmail.com}{fr.vicidomini94@gmail.com}\par
      \href{https://francescovicidomini.dev/}{francescovicidomini.dev}\par
      \href{https://github.com/CiccioTecchio}{github.com/CiccioTecchio}
    }
    \vfill
  \end{frame}
}

\end{document}
